\section{Теория}

Вспомним, что такое \textbf{мощность множества}. Мы будем работать только с конечными множествами. Мощность конечного множества — это количество его элементов.

\begin{designation}
    Мощность множества \(A\) обозначается \(\lvert A\rvert\).
\end{designation}

\subsection{Правило сложения}

\begin{example}[combinatorics-rules/theory/1]\label{example-theory-1}
\textbf{Задача.} Вам задали домашнее задание по алгебре, английскому, русскому и теории вероятностей. По алгебре нужно сделать 5 упражнений, по английскому — 8, по русскому — 6, по теории вероятностей — 3. Сколько существует способов выбрать одно упражнение?
\end{example}

\begin{solution}[combinatorics-rules/theory/1]\label{solution-combinatorics-rules-theory-1}
Соберём все упражнения в один список: 5 по алгебре, 8 по английскому, 6 по русскому и 3 по теории вероятностей.

Всего получится $5+8+6+3=22$ упражнения. Значит, выбрать одно из них можно 22 способами.

Мы просто сложили четыре числа. Это сработало, потому что ни одно упражнение не попало в список дважды. Так мы приходим к правилу сложения.
\end{solution}

\begin{definition}
    \textbf{Правило сложения.} Пусть \(A_1,A_2,\dots,A_n\) — попарно непересекающиеся множества. Количество способов выбрать один элемент из их объединения равно сумме мощностей этих множеств.
    \[
        \lvert A_1\cup A_2\cup\dots\cup A_n\rvert=\lvert A_1\rvert+\lvert A_2\rvert+\dots+\lvert A_n\rvert.
    \]
\end{definition}

\subsection{Правило умножения}

\begin{definition}
    Пусть даны множества \(A\) и \(B\). Возьмём элемент \(a\) из \(A\) и элемент \(b\) из \(B\). Запись \((a,b)\) называется \textbf{упорядоченной парой} элементов из множеств \(A\) и \(B\). На первом месте стоит элемент из \(A\), на втором — элемент из \(B\).
\end{definition}

Для большего числа множеств так же определяются упорядоченные наборы.

В наборе \((a_1,a_2,\dots,a_n)\) на \(i\)-м месте стоит элемент \(a_i\) из множества \(A_i\).

\begin{example}[combinatorics-rules/theory/2]\label{example-theory-2}
\textbf{Задача.} Вам задали то же домашнее задание. Сколько существует способов выбрать по одному упражнению по каждому предмету?
\end{example}

\begin{solution}[combinatorics-rules/theory/2]\label{solution-combinatorics-rules-theory-2}
Для каждого из 5 упражнений по алгебре можно выбрать любое из 8 упражнений по английскому. Получится $5\cdot8=40$ пар.

К каждой из 40 пар можно добавить любое из 6 упражнений по русскому. Получится $40\cdot6=240$ троек.

К каждой из 240 троек можно добавить любое из 3 упражнений по теории вероятностей. Получится $240\cdot3=720$ вариантов.

\textbf{Ответ:} 720 способов. На каждом шаге мы умножали число уже полученных вариантов на число новых вариантов. В этом и состоит идея правила умножения.
\end{solution}

\begin{definition}
    \textbf{Правило умножения.} Количество упорядоченных наборов \((a_1,a_2,\dots,a_n)\), где \(a_i\in A_i\) для каждого \(i\), равно произведению мощностей этих множеств.
    \[
        \lvert A_1\rvert\cdot\lvert A_2\rvert\cdot\dots\cdot\lvert A_n\rvert.
    \]
\end{definition}

\section{Задачи}

\begin{task}[combinatorics-rules/problems/1]\label{task-problems-1}
\textbf{1.} Сколько на данный момент может быть зарегистрировано легковых машин в Москве? Напомню, что в Москве используются коды региона: 77, 97, 99, 177, 199, 197, 777, 797, 799, 977, 277, 299. Буквы в номерном знаке автомобиля используются только имеющие аналоги в латинице, то есть А, В, Е, К, М, Н, О, Р, С, Т, У, Х.
\end{task}

\begin{solution}[combinatorics-rules/problems/1]\label{solution-combinatorics-rules-problems-1}
В номере три буквы. На каждое место можно поставить любую из 12 букв, поэтому сочетаний букв будет $12\cdot12\cdot12$.

Для каждого сочетания букв можно выбрать любое число от 001 до 999 — это 999 вариантов. Ещё можно выбрать один из 12 кодов региона из условия.

Всего получаем
\[
12\cdot12\cdot12\cdot999\cdot12=20\,715\,264.
\]
\textbf{Ответ:} $20\,715\,264$ номера. Здесь считаем, что все перечисленные сочетания можно использовать.
\end{solution}

\begin{task}[combinatorics-rules/problems/2]\label{task-problems-2}
\textbf{2.} Сколько рёбер в полном графе на \(n\) вершинах?
\end{task}

\begin{solution}[combinatorics-rules/problems/2]\label{solution-combinatorics-rules-problems-2}
\textbf{Через правило умножения.} Каждая вершина соединена с остальными $n-1$ вершинами.

Вершин $n$, поэтому сначала получаем $n(n-1)$. Но каждое ребро мы посчитали дважды — с одного конца и с другого. Делим на 2:
\[
\frac{n(n-1)}{2}.
\]
\end{solution}

\begin{solution}[combinatorics-rules/problems/2]\label{solution-combinatorics-rules-problems-2-addition}
\textbf{Через правило сложения.} Пронумеруем вершины от 1 до $n$.

Сначала посчитаем $n-1$ рёбер, выходящих из первой вершины. У второй осталось $n-2$ ещё не посчитанных ребра: ребро к первой уже учтено. У третьей останется $n-3$, и так далее. У последней новых рёбер не останется.

Значит, нужно сложить числа от 0 до $n-1$. Обозначим сумму буквой $S$ и запишем её дважды, во второй раз — в обратном порядке:
\[
\begin{aligned}
S&=0+1+\dots+(n-2)+(n-1),\\
S&=(n-1)+(n-2)+\dots+1+0.
\end{aligned}
\]
Складываем числа друг под другом. В каждом столбце получается $n-1$, а столбцов $n$. Поэтому $2S=n(n-1)$ и
\[
S=\frac{n(n-1)}{2}.
\]
\end{solution}

\begin{task}[combinatorics-rules/problems/3]\label{task-problems-3}
\textbf{3.} На дипломатической встрече было 5 делегаций. В первой 7 человек, во второй 5, в третьей 8, в четвёртой 7, в пятой 3. Каждый дипломат жмёт руку каждому иностранному дипломату (то есть жмёт всем, кроме соотечественников). Сколько всего рукопожатий было сделано?
\end{task}

\begin{solution}[combinatorics-rules/problems/3]\label{solution-combinatorics-rules-problems-3}
Будем считать рукопожатия так, чтобы ни одно не повторилось.

Каждый из 7 человек первой делегации пожмёт руки $5+8+7+3=23$ людям из остальных делегаций. Это $7\cdot23=161$ рукопожатие.

Для второй делегации считаем только встречи с третьей, четвёртой и пятой: встречи с первой уже учтены. Получаем $5\cdot(8+7+3)=90$ рукопожатий.

Для третьей остались встречи с четвёртой и пятой: $8\cdot(7+3)=80$.

Для четвёртой остались только встречи с пятой: $7\cdot3=21$.

Складываем:
\[
161+90+80+21=352.
\]
\textbf{Ответ:} 352 рукопожатия.
\end{solution}

\section{Домашнее задание}

\begin{task}[combinatorics-rules/homework/1]\label{task-homework-1}
\textbf{1.} Наугад выбирают машину, зарегистрированную в Москве. Какова вероятность того, что её номер будет содержать буквы АМР именно в таком порядке? Ответ округлите до 3 знаков после запятой.
\end{task}

\begin{solution}[combinatorics-rules/homework/1]\label{solution-combinatorics-rules-homework-1}
\textbf{1. Сколько всего номеров.} На уроке мы получили $20\,715\,264$ номера.

\textbf{2. Сколько номеров с буквами АМР.} Буквы уже выбраны. Остаётся 999 вариантов числа и 12 кодов региона. Значит, таких номеров $999\cdot12=11\,988$.

\textbf{3. Найдём вероятность.} Считаем все эти номера равновероятными. Делим число подходящих номеров на общее число номеров:
\[
P=\frac{11\,988}{20\,715\,264}\approx0{,}000579.
\]
\textbf{Ответ:} $0{,}001$ после округления до трёх знаков после запятой.
\end{solution}

\begin{task}[combinatorics-rules/homework/2]\label{task-homework-2}
\textbf{2.} (№ 141) В школе есть все классы — с 1 по 11. Каждый из них имеет дополнительную букву А, Б, В, Г или Д. Например, класс 1А, 7Б и т. д. Сколько всего классов в этой школе?
\end{task}

\begin{solution}[combinatorics-rules/homework/2]\label{solution-combinatorics-rules-homework-2}
Считаем, что в каждой параллели есть все пять классов: А, Б, В, Г и Д.

В первой параллели 5 классов, во второй ещё 5, и так до одиннадцатой. Всего $11\cdot5=55$ классов.

\textbf{Ответ:} 55 классов.
\end{solution}

\begin{task}[combinatorics-rules/homework/3]\label{task-homework-3}
\textbf{3.} (№ 143) На день школы каждая девочка подарила каждому мальчику открытку, а каждый мальчик подарил каждой девочке гвоздику. Чего было больше — подаренных открыток или подаренных гвоздик?
\end{task}

\begin{solution}[combinatorics-rules/homework/3]\label{solution-combinatorics-rules-homework-3}
Возьмём любую девочку и любого мальчика. Девочка дарит ему одну открытку, а мальчик ей — одну гвоздику.

Так происходит с каждой парой «девочка — мальчик»: на одну открытку приходится одна гвоздика. Значит, открыток и гвоздик подарили поровну.
\end{solution}

\begin{task}[combinatorics-rules/homework/4]\label{task-homework-4}
\textbf{4.} (№ 144) Составляются различные последовательности из цифр 0 и 1 (\textbf{двоичные последовательности}). Сколько существует двоичных последовательностей длины:\\
а) 1;    б) 3;    в) 10;    г) \(n\)?
\end{task}

\begin{solution}[combinatorics-rules/homework/4]\label{solution-combinatorics-rules-homework-4}
На каждом месте можно написать 0 или 1 — два варианта. Когда добавляем ещё одно место, число последовательностей удваивается.

Поэтому для длины $n$ нужно перемножить $n$ двоек: получится $2^n$.

\textbf{Ответ:} а) $2$; б) $2^3=8$; в) $2^{10}=1024$; г) $2^n$.
\end{solution}

\begin{task}[combinatorics-rules/homework/5]\label{task-homework-5}
\textbf{5.} (№ 145) На приёме в посольстве встретились две делегации, в каждой из которых было несколько дипломатов (больше одного). Каждый дипломат одной делегации пожал руку каждому дипломату второй делегации. Сколько было членов в каждой делегации, если всего произошло 143 рукопожатия?
\end{task}

\begin{solution}[combinatorics-rules/homework/5]\label{solution-combinatorics-rules-homework-5}
Если в первой делегации $a$ человек, а во второй $b$, то каждый из $a$ человек пожмёт руки всем $b$ гостям. Всего получится $a\cdot b=143$ рукопожатия.

Разложим 143 на множители: $143=11\cdot13$. Числа 11 и 13 простые. Другой вариант — $1\cdot143$, но он не подходит: в каждой делегации больше одного человека.

\textbf{Ответ:} 11 человек и 13 человек.
\end{solution}

\begin{task}[combinatorics-rules/homework/6]\label{task-homework-6}
\textbf{6.} (№ 146) Важные данные в компьютере часто защищают паролем. Для пароля можно использовать большие и малые латинские буквы, цифры, некоторые знаки. Всего разрешённых символов 92. Составьте числовое выражение для общего числа возможных паролей, состоящих из 8 символов. Пользуясь калькулятором, вычислите приближённо его значение.
\end{task}

\begin{solution}[combinatorics-rules/homework/6]\label{solution-combinatorics-rules-homework-6}
На первое место можно поставить любой из 92 символов. На второе тоже 92, и так на каждое из восьми мест.

Значит, нужно перемножить восемь чисел 92:
\[
92^8=5\,132\,188\,731\,375\,616\approx5{,}13\cdot10^{15}.
\]
\end{solution}
